\documentclass[12pt,a4paper]{article}
\usepackage[utf8]{inputenc}
\usepackage[T2A]{fontenc}
\usepackage[russian]{babel}
\usepackage{amsmath,amssymb,amsthm}
\usepackage{geometry}
\geometry{left=25mm, right=25mm, top=20mm, bottom=20mm}

\title{Метод интегральных соотношений Баренблатта\\
для задачи двухфазной фильтрации\\
в модели Black Oil}
\author{}
\date{}

\begin{document}
\maketitle

%=============================================================
\section{Постановка задачи}
%=============================================================

Рассматривается одномерная изотермическая фильтрация нефти с растворённым газом
в пористой среде. Основные допущения:
\begin{itemize}
  \item две фазы: нефть~(o) и газ~(g); воды нет;
  \item начальное давление $p_0 > p_b$ (выше давления насыщения);
  \item пористая среда недеформируема, пористость $\phi = \mathrm{const}$;
  \item гравитационные эффекты не учитываются;
  \item термодинамическое равновесие между фазами (мгновенный массообмен).
\end{itemize}

По мере снижения давления ниже~$p_b$ газ выделяется из нефти
в свободную фазу. Скорость фильтрации каждой фазы определяется
обобщённым законом Дарси с относительными фазовыми проницаемостями.

%=============================================================
\section{Определяющие соотношения модели Black Oil}
%=============================================================

\subsection{Насыщенности}

\begin{equation}\label{eq:saturation}
  S_o + S_g = 1.
\end{equation}

\subsection{PVT-свойства}

Модель Black Oil задаёт следующие зависимости от давления~$p$:
\begin{itemize}
  \item $B_o(p)$ --- объёмный коэффициент нефти (пластовый объём / стандартный объём);
  \item $B_g(p)$ --- объёмный коэффициент газа;
  \item $R_s(p)$ --- газосодержание нефти (объём растворённого газа при стандартных
        условиях на единицу стандартного объёма нефти);
  \item $\mu_o(p)$, $\mu_g(p)$ --- вязкости нефти и газа.
\end{itemize}

При $p \ge p_b$ газосодержание постоянно: $R_s = R_s(p_b) \equiv R_{si}$,
свободного газа нет ($S_g = 0$).

\subsection{Законы сохранения массы}

\textbf{Нефтяной компонент:}
\begin{equation}\label{eq:oil_conservation}
  \frac{\partial}{\partial t}\left(\phi\,\frac{S_o}{B_o}\right)
  + \frac{\partial}{\partial x}\left(\frac{u_o}{B_o}\right) = 0.
\end{equation}

\textbf{Газовый компонент} (свободный газ + газ, растворённый в нефти):
\begin{equation}\label{eq:gas_conservation}
  \frac{\partial}{\partial t}\left(\phi\left(\frac{S_g}{B_g}
  + \frac{R_s\,S_o}{B_o}\right)\right)
  + \frac{\partial}{\partial x}\left(\frac{u_g}{B_g}
  + \frac{R_s\,u_o}{B_o}\right) = 0.
\end{equation}

\subsection{Закон Дарси}

\begin{equation}\label{eq:darcy_oil}
  u_o = -\frac{k\,k_{ro}(S_g)}{\mu_o(p)}\,\frac{\partial p}{\partial x},
\end{equation}

\begin{equation}\label{eq:darcy_gas}
  u_g = -\frac{k\,k_{rg}(S_g)}{\mu_g(p)}\,\frac{\partial p}{\partial x},
\end{equation}
где $k$ --- абсолютная проницаемость, $k_{ro}(S_g)$ и $k_{rg}(S_g)$ ---
относительные фазовые проницаемости.

\subsection{Суммарная скорость и дробный поток}

Суммарная скорость фильтрации:
\begin{equation}\label{eq:total_velocity}
  u = u_o + u_g = -k\,\lambda_t(S_g,\,p)\,\frac{\partial p}{\partial x},
\end{equation}
где суммарная подвижность:
\begin{equation}\label{eq:total_mobility}
  \lambda_t = \frac{k_{ro}}{\mu_o} + \frac{k_{rg}}{\mu_g}.
\end{equation}

Дробный поток нефти:
\begin{equation}\label{eq:fractional_flow}
  f_o = \frac{\lambda_o}{\lambda_t}
  = \frac{k_{ro}/\mu_o}{k_{ro}/\mu_o + k_{rg}/\mu_g}.
\end{equation}

%=============================================================
\section{Структура решения: две зоны}
%=============================================================

По мере снижения давления на забое ниже~$p_b$ возникают две
характерные области (рис.~\ref{fig:zones}):

\begin{description}
  \item[Зона~I] ($0 \le x \le x_f(t)$): давление $p < p_b$, двухфазная
    фильтрация, газонасыщенность $S_g > 0$.
  \item[Зона~II] ($x_f(t) < x < x_e(t)$): давление $p \ge p_b$, однофазная
    фильтрация нефти, $S_g = 0$.
\end{description}

За пределами фронта возмущения ($x > x_e(t)$) пласт не возмущён: $p = p_0$.

На \textbf{фронте насыщения} $x = x_f(t)$:
\begin{equation}\label{eq:front_conditions}
  p\big|_{x=x_f} = p_b, \qquad S_g\big|_{x=x_f} = 0.
\end{equation}

Схематически:
\[
  \underbrace{\text{Забой}
  \;\xrightarrow{\;S_g > 0,\; p < p_b\;}}_{\text{Зона I}}
  \; x_f(t) \;
  \underbrace{\xrightarrow{\;S_g = 0,\; p \ge p_b\;}}_{\text{Зона II}}
  \; x_e(t) \;
  \xrightarrow{\;p = p_0\;}
  \text{Невозмущённый пласт}
\]

\begin{figure}[h]
  \centering
  \setlength{\unitlength}{1mm}
  \begin{picture}(120,40)
    \put(0,5){\vector(1,0){120}}
    \put(122,4){$x$}
    \put(0,5){\vector(0,1){35}}
    \put(-5,42){$p, S_g$}
    % pressure profile
    \qbezier(0,15)(20,22)(50,30)
    \qbezier(50,30)(70,35)(90,37)
    \put(90,37){\line(1,0){25}}
    % saturation profile
    \qbezier(0,25)(15,18)(35,10)
    \qbezier(35,10)(45,7)(50,5)
    % labels
    \put(50,1){\line(0,1){4}}
    \put(47,-2){$x_f$}
    \put(90,1){\line(0,1){4}}
    \put(87,-2){$x_e$}
    \put(5,11){$p(x)$}
    \put(5,27){$S_g(x)$}
    \put(115,36){$p_0$}
    \put(52,29){$p_b$}
  \end{picture}
  \caption{Схема распределения давления и газонасыщенности.}
  \label{fig:zones}
\end{figure}

%=============================================================
\section{Метод интегральных соотношений}
%=============================================================

\subsection{Идея метода}

Следуя Баренблатту, вместо решения системы уравнений в частных
производных~\eqref{eq:oil_conservation}--\eqref{eq:gas_conservation}
мы:
\begin{enumerate}
  \item интегрируем уравнения сохранения по пространственной переменной
        в пределах каждой зоны;
  \item аппроксимируем профили давления и насыщенности функциями
        заданной формы с неизвестными параметрами;
  \item подставляем аппроксимации в интегральные соотношения и получаем
        систему ОДУ для параметров.
\end{enumerate}

\subsection{Интегрирование уравнения сохранения нефти (Зона~I)}

Интегрируем~\eqref{eq:oil_conservation} по~$x$ от~$0$ до~$x_f(t)$:
\begin{equation}\label{eq:oil_integral_raw}
  \int_0^{x_f} \frac{\partial}{\partial t}
  \left(\phi\,\frac{S_o}{B_o}\right) dx
  + \left.\frac{u_o}{B_o}\right|_0^{x_f} = 0.
\end{equation}

Применяя формулу Лейбница к первому слагаемому:
\begin{equation}\label{eq:leibniz}
  \int_0^{x_f}\frac{\partial}{\partial t}(\cdot)\,dx
  = \frac{d}{dt}\int_0^{x_f}(\cdot)\,dx
  - (\cdot)\big|_{x=x_f}\,\dot{x}_f,
\end{equation}
получаем с учётом $S_o(x_f) = 1$, $p(x_f) = p_b$:
\begin{equation}\label{eq:oil_integral}
  \boxed{
  \frac{d}{dt}\int_0^{x_f} \phi\,\frac{1 - S_g}{B_o(p)}\,dx
  = \frac{\phi}{B_o(p_b)}\,\dot{x}_f
  - \frac{u_o(x_f)}{B_o(p_b)}
  + \frac{u_o(0)}{B_o(p_{wf})}.
  }
\end{equation}

\subsection{Интегрирование уравнения сохранения газа (Зона~I)}

Аналогично интегрируем~\eqref{eq:gas_conservation}:
\begin{equation}\label{eq:gas_integral}
  \boxed{
  \frac{d}{dt}\int_0^{x_f} \phi\left(\frac{S_g}{B_g}
  + \frac{R_s(1-S_g)}{B_o}\right)dx
  = \phi\,\frac{R_s(p_b)}{B_o(p_b)}\,\dot{x}_f
  - \left.\left(\frac{u_g}{B_g}+\frac{R_s\,u_o}{B_o}\right)\right|_0^{x_f}.
  }
\end{equation}

На фронте $x = x_f$ свободного газа нет ($S_g = 0$, $k_{rg} = 0$), поэтому
$u_g(x_f) = 0$, и правая часть упрощается.

%=============================================================
\section{Аппроксимация профилей}
%=============================================================

Введём автомодельную переменную:
\begin{equation}\label{eq:xi}
  \xi = \frac{x}{x_f(t)}, \qquad 0 \le \xi \le 1.
\end{equation}

\subsection{Профиль давления в Зоне~I}

\begin{equation}\label{eq:pressure_profile}
  p(\xi, t) = p_b - \Delta p(t)\,(1-\xi)^n,
\end{equation}
где $\Delta p(t) = p_b - p_{wf}(t)$ --- депрессия на забое,
$n > 0$ --- показатель формы.

Граничные условия:
\[
  p(0,t) = p_b - \Delta p = p_{wf}(t), \qquad p(1,t) = p_b.
\]

Градиент давления:
\begin{equation}\label{eq:dp_dx}
  \frac{\partial p}{\partial x}
  = \frac{n\,\Delta p}{x_f}\,(1-\xi)^{n-1}.
\end{equation}

\subsection{Профиль газонасыщенности в Зоне~I}

\begin{equation}\label{eq:saturation_profile}
  S_g(\xi, t) = S_{g,\max}(t)\,(1-\xi)^m,
\end{equation}
где $S_{g,\max}(t)$ --- максимальная газонасыщенность (на забое, $\xi=0$),
$m > 0$ --- показатель формы.

Граничные условия:
\[
  S_g(0,t) = S_{g,\max}(t), \qquad S_g(1,t) = 0.
\]

\subsection{Вычисление интегралов}

Благодаря степенной форме профилей интегралы вычисляются аналитически.

\textbf{Интеграл газонасыщенности:}
\begin{equation}\label{eq:Sg_integral}
  \int_0^{x_f} S_g\,dx
  = S_{g,\max}\,x_f \int_0^1 (1-\xi)^m\,d\xi
  = \frac{S_{g,\max}\,x_f}{m+1}.
\end{equation}

\textbf{Интеграл нефтенасыщенности:}
\begin{equation}\label{eq:So_integral}
  \int_0^{x_f}(1-S_g)\,dx
  = x_f\left(1 - \frac{S_{g,\max}}{m+1}\right).
\end{equation}

Для более точного вычисления интегралов, содержащих PVT-свойства,
линеаризуем их вблизи~$p_b$:
\begin{equation}\label{eq:Bo_linearize}
  \frac{1}{B_o(p)} \approx \frac{1}{B_o(p_b)}
  \left(1 + c_o\,(p_b - p)\right),
\end{equation}
где $c_o = -\frac{1}{B_o}\frac{dB_o}{dp}\big|_{p_b}$ ---
сжимаемость нефти (с учётом выделения газа).

Тогда:
\begin{equation}\label{eq:oil_volume_integral}
  \int_0^{x_f}\frac{1-S_g}{B_o}\,dx
  \approx \frac{x_f}{B_o(p_b)}\left[
    1 - \frac{S_{g,\max}}{m+1}
    + c_o\,\Delta p\left(\frac{1}{n+1}
    - \frac{S_{g,\max}}{m+n+1}\right)
  \right].
\end{equation}

\subsection{Профиль давления в Зоне~II}

В однофазной области ($x_f < x < x_e$) вводим переменную:
\begin{equation}\label{eq:eta}
  \eta = \frac{x - x_f(t)}{x_e(t) - x_f(t)}, \qquad 0 \le \eta \le 1,
\end{equation}
и аппроксимируем:
\begin{equation}\label{eq:pressure_zone2}
  p(\eta,t) = p_b + (p_0 - p_b)\,\eta^{n_2},
\end{equation}
так что $p = p_b$ при $\eta = 0$ (внутренняя граница)
и $p = p_0$ при $\eta = 1$ (фронт возмущения).

Интегральное соотношение для Зоны~II получается интегрированием
уравнения пьезопроводности:
\begin{equation}\label{eq:zone2_integral}
  \frac{d}{dt}\int_{x_f}^{x_e}\phi\,c_t\,(p - p_0)\,dx
  = \left.\frac{k}{\mu_o\,B_o}\frac{\partial p}{\partial x}
  \right|_{x_f}
  - \left.\frac{k}{\mu_o\,B_o}\frac{\partial p}{\partial x}
  \right|_{x_e},
\end{equation}
где последний член обращается в нуль, если на фронте возмущения
$\partial p/\partial x = 0$.

%=============================================================
\section{Замкнутая система ОДУ}
%=============================================================

После подстановки аппроксимаций в интегральные
соотношения~\eqref{eq:oil_integral}--\eqref{eq:gas_integral}
и~\eqref{eq:zone2_integral} получаем систему обыкновенных
дифференциальных уравнений для четырёх неизвестных функций времени:

\begin{equation}\label{eq:unknowns}
  x_f(t), \quad x_e(t), \quad \Delta p(t), \quad S_{g,\max}(t).
\end{equation}

\subsection{Уравнение из сохранения нефти}

Обозначим для краткости:
\[
  \mathcal{I}_o = \frac{\phi\,x_f}{B_o(p_b)}
  \left[1 - \frac{S_{g,\max}}{m+1}
  + c_o\,\Delta p\left(\frac{1}{n+1}
  - \frac{S_{g,\max}}{m+n+1}\right)\right].
\]

Тогда~\eqref{eq:oil_integral} принимает вид:
\begin{equation}\label{eq:ode_oil}
  \frac{d\mathcal{I}_o}{dt}
  = \frac{\phi}{B_o(p_b)}\,\dot{x}_f
  + \frac{k}{\mu_o(p_b)\,B_o(p_b)}\,
  \frac{n\,\Delta p}{x_f}\,k_{ro}(0)
  - \frac{k}{\mu_o(p_{wf})\,B_o(p_{wf})}\,
  \frac{n\,\Delta p}{x_f}\,k_{ro}(S_{g,\max}).
\end{equation}

Здесь использовано~\eqref{eq:dp_dx} для вычисления скоростей на границах.
На фронте ($\xi=1$): $k_{ro}(S_g=0) = k_{ro}(0) \equiv 1$.
На забое ($\xi=0$): $k_{ro}(S_{g,\max})$.

\subsection{Уравнение из сохранения газа}

Аналогично, определив:
\[
  \mathcal{I}_g = \phi\,x_f\left[
    \frac{S_{g,\max}}{(m+1)\,B_g^*}
    + \frac{R_s^*}{B_o^*}\left(1 - \frac{S_{g,\max}}{m+1}\right)
  \right],
\]
где звёздочкой обозначены средневзвешенные по зоне значения PVT-свойств,
получаем из~\eqref{eq:gas_integral}:
\begin{equation}\label{eq:ode_gas}
  \frac{d\mathcal{I}_g}{dt}
  = \frac{\phi\,R_s(p_b)}{B_o(p_b)}\,\dot{x}_f
  - \frac{R_s(p_b)}{B_o(p_b)}\,u_o(x_f)
  + \frac{1}{B_g(p_{wf})}\,u_g(0)
  + \frac{R_s(p_{wf})}{B_o(p_{wf})}\,u_o(0).
\end{equation}

\subsection{Уравнение для Зоны~II}

Из~\eqref{eq:zone2_integral}:
\begin{equation}\label{eq:ode_zone2}
  \frac{d}{dt}\left[
    \frac{\phi\,c_t\,(p_0 - p_b)\,(x_e - x_f)}{n_2 + 1}
  \right]
  = \frac{k}{\mu_o\,B_o}\,
  \frac{n_2\,(p_0 - p_b)}{x_e - x_f}.
\end{equation}

\subsection{Условие сшивки потоков}

Четвёртое уравнение --- непрерывность суммарного потока на фронте~$x_f$:
\begin{equation}\label{eq:flux_matching}
  u\big|_{x_f^-} = u\big|_{x_f^+},
\end{equation}
что связывает градиенты давления по обе стороны фронта:
\begin{equation}\label{eq:gradient_matching}
  \frac{k\,\lambda_t(0,\,p_b)}{1}\cdot\frac{n\,\Delta p}{x_f}
  = \frac{k}{\mu_o(p_b)\,B_o(p_b)}\cdot
  \frac{n_2\,(p_0 - p_b)}{x_e - x_f}.
\end{equation}

Здесь слева: двухфазная подвижность при $S_g=0$ на фронте,
справа: однофазная подвижность в Зоне~II.

\subsection{Граничное условие}

Система замыкается заданием условия на забое. Например:
\begin{itemize}
  \item постоянный дебит нефти:
    $\displaystyle q_o = -\frac{k\,k_{ro}(S_{g,\max})}{\mu_o(p_{wf})}\,
    \frac{n\,\Delta p}{x_f}\bigg|_{x=0} = \mathrm{const}$;
  \item постоянное забойное давление: $p_{wf} = \mathrm{const}$,
    т.\,е. $\Delta p = p_b - p_{wf} = \mathrm{const}$.
\end{itemize}

%=============================================================
\section{Решение системы ОДУ}
%=============================================================

Итоговая система~\eqref{eq:ode_oil}--\eqref{eq:gradient_matching}
с граничным условием представляет собой начальную задачу для
вектора $\{x_f,\, x_e,\, \Delta p,\, S_{g,\max}\}$.

Начальные условия соответствуют моменту, когда забойное давление
достигает~$p_b$:
\[
  x_f(0) = 0, \quad x_e(0) = x_e^{(0)}, \quad
  \Delta p(0) = 0, \quad S_{g,\max}(0) = 0.
\]

Система решается численно методом Рунге--Кутты (например, RK4).

\subsection{Результаты решения}

Решение даёт:
\begin{enumerate}
  \item \textbf{Динамику фронта двухфазной зоны} $x_f(t)$ ---
    скорость продвижения области с выделившимся газом.
  \item \textbf{Эволюцию газонасыщенности} $S_{g,\max}(t)$ ---
    накопление свободного газа в призабойной зоне.
  \item \textbf{Динамику давления} $\Delta p(t)$ и $p_{wf}(t)$ ---
    изменение забойного давления и депрессии.
  \item \textbf{Газовый фактор} --- отношение дебитов газа и нефти
    на забое, определяемое через дробный поток.
\end{enumerate}

%=============================================================
\section{Радиальная геометрия}
%=============================================================

Для радиальной постановки (скважина радиуса~$r_w$) все уравнения
переписываются в цилиндрических координатах. Уравнение сохранения
нефти:
\begin{equation}\label{eq:oil_radial}
  \frac{\partial}{\partial t}\left(\phi\,\frac{S_o}{B_o}\right)
  + \frac{1}{r}\frac{\partial}{\partial r}
  \left(r\,\frac{u_o}{B_o}\right) = 0,
\end{equation}
а интегрирование ведётся по~$r$ от~$r_w$ до~$r_f(t)$:
\begin{equation}\label{eq:oil_integral_radial}
  \frac{d}{dt}\int_{r_w}^{r_f}
  \phi\,\frac{1-S_g}{B_o}\,r\,dr
  = \frac{\phi\,r_f}{B_o(p_b)}\,\dot{r}_f
  - \frac{r_f\,u_o(r_f)}{B_o(p_b)}
  + \frac{r_w\,u_o(r_w)}{B_o(p_{wf})}.
\end{equation}

Аппроксимация профилей в этом случае строится через переменную:
\begin{equation}
  \xi = \frac{\ln(r/r_w)}{\ln(r_f/r_w)}, \qquad 0 \le \xi \le 1,
\end{equation}
что учитывает логарифмический характер распределения давления
вблизи скважины.

%=============================================================
\section{Определение показателей формы методом моментов}
%=============================================================

В предыдущих разделах показатели формы $n$, $m$ (и $n_2$) были
свободными параметрами. Метод моментов позволяет определить их
из дополнительных интегральных соотношений, так что они становятся
частью решения.

\subsection{Идея метода}

Исходное уравнение в частных производных содержит бесконечное
количество информации о профиле решения. Интегрирование
$\int_0^{x_f}(\cdot)\,dx$ извлекает лишь \emph{нулевой момент} ---
суммарное количество вещества в зоне. Умножая уравнение на степени
координаты перед интегрированием, получаем \emph{моменты высших
порядков}, которые несут информацию о \emph{форме} профиля.

Каждый дополнительный момент даёт одно дополнительное уравнение,
позволяющее найти один дополнительный параметр.

\subsection{Первый момент уравнения сохранения нефти}

Умножим~\eqref{eq:oil_conservation} на~$x$ и проинтегрируем
по~$x$ от~$0$ до~$x_f$:
\begin{equation}\label{eq:first_moment_oil}
  \int_0^{x_f} x\,\frac{\partial}{\partial t}
  \left(\phi\,\frac{S_o}{B_o}\right)dx
  + \int_0^{x_f} x\,\frac{\partial}{\partial x}
  \left(\frac{u_o}{B_o}\right)dx = 0.
\end{equation}

Первый интеграл преобразуем по формуле Лейбница:
\begin{equation}\label{eq:first_moment_leibniz}
  \int_0^{x_f} x\,\frac{\partial}{\partial t}
  \left(\phi\,\frac{S_o}{B_o}\right)dx
  = \frac{d}{dt}\int_0^{x_f} x\,\phi\,\frac{S_o}{B_o}\,dx
  - x_f\,\frac{\phi}{B_o(p_b)}\,\dot{x}_f.
\end{equation}

Второй интеграл берём по частям:
\begin{equation}\label{eq:first_moment_parts}
  \int_0^{x_f} x\,\frac{\partial}{\partial x}
  \left(\frac{u_o}{B_o}\right)dx
  = \left.x\,\frac{u_o}{B_o}\right|_0^{x_f}
  - \int_0^{x_f}\frac{u_o}{B_o}\,dx
  = x_f\,\frac{u_o(x_f)}{B_o(p_b)}
  - \int_0^{x_f}\frac{u_o}{B_o}\,dx.
\end{equation}

Объединяя, получаем \textbf{первый момент}:
\begin{equation}\label{eq:moment1_oil}
  \boxed{
  \frac{d}{dt}\int_0^{x_f} x\,\phi\,\frac{1-S_g}{B_o}\,dx
  = x_f\left(\frac{\phi}{B_o(p_b)}\,\dot{x}_f
  - \frac{u_o(x_f)}{B_o(p_b)}\right)
  + \int_0^{x_f}\frac{u_o}{B_o}\,dx.
  }
\end{equation}

\subsection{Вычисление моментных интегралов}

Подставляя аппроксимации~\eqref{eq:pressure_profile}
и~\eqref{eq:saturation_profile}, получаем (в ведущем порядке
по~$c_o\,\Delta p$):

\textbf{Первый момент нефтенасыщенности:}
\begin{equation}\label{eq:x_So_integral}
  \int_0^{x_f} x\,(1-S_g)\,dx
  = x_f^2\int_0^1 \xi\,(1 - S_{g,\max}(1-\xi)^m)\,d\xi
  = x_f^2\left[\frac{1}{2}
  - S_{g,\max}\left(\frac{1}{m+1} - \frac{1}{m+2}\right)\right].
\end{equation}

Упрощая:
\begin{equation}\label{eq:x_So_integral_simple}
  \int_0^{x_f} x\,(1-S_g)\,dx
  = x_f^2\left[\frac{1}{2}
  - \frac{S_{g,\max}}{(m+1)(m+2)}\right].
\end{equation}

\textbf{Интеграл потока нефти:}
\begin{equation}\label{eq:uo_Bo_integral}
  \int_0^{x_f}\frac{u_o}{B_o}\,dx
  = -\frac{k\,n\,\Delta p}{x_f}
  \int_0^1 \frac{k_{ro}(S_g(\xi))}
  {\mu_o(p(\xi))\,B_o(p(\xi))}\,(1-\xi)^{n-1}\,d\xi.
\end{equation}

Этот интеграл в общем случае вычисляется численно (квадратурой),
но при малых~$S_g$ допускает линеаризацию:
\begin{equation}\label{eq:kro_linearize}
  k_{ro}(S_g) \approx 1 - \alpha\,S_g,
  \qquad \alpha = -k_{ro}'(0),
\end{equation}
что позволяет выразить его аналитически через бета-функции.

\subsection{Первый момент уравнения сохранения газа}

Аналогично умножаем~\eqref{eq:gas_conservation} на~$x$
и интегрируем:
\begin{equation}\label{eq:moment1_gas}
  \boxed{
  \frac{d}{dt}\int_0^{x_f} x\,\phi
  \left(\frac{S_g}{B_g}+\frac{R_s(1-S_g)}{B_o}\right)dx
  = x_f\left(\frac{\phi\,R_s(p_b)}{B_o(p_b)}\,\dot{x}_f
  - \frac{R_s(p_b)\,u_o(x_f)}{B_o(p_b)}\right)
  + \int_0^{x_f}\left(\frac{u_g}{B_g}
  + \frac{R_s\,u_o}{B_o}\right)dx.
  }
\end{equation}

\subsection{Расширенная система ОДУ}

С учётом моментных соотношений~\eqref{eq:moment1_oil}
и~\eqref{eq:moment1_gas} система неизвестных расширяется:
\begin{equation}\label{eq:unknowns_extended}
  x_f(t), \quad x_e(t), \quad \Delta p(t), \quad
  S_{g,\max}(t), \quad n(t), \quad m(t).
\end{equation}

Теперь имеем шесть уравнений для шести неизвестных:
\begin{enumerate}
  \item нулевой момент сохранения нефти~\eqref{eq:oil_integral}
        $\longrightarrow$ связывает $x_f$, $\Delta p$, $S_{g,\max}$;
  \item нулевой момент сохранения газа~\eqref{eq:gas_integral}
        $\longrightarrow$ связывает $x_f$, $\Delta p$, $S_{g,\max}$;
  \item интегральное соотношение для Зоны~II~\eqref{eq:zone2_integral}
        $\longrightarrow$ определяет $x_e$;
  \item сшивка потоков~\eqref{eq:flux_matching}
        $\longrightarrow$ алгебраическая связь;
  \item первый момент сохранения нефти~\eqref{eq:moment1_oil}
        $\longrightarrow$ определяет $n(t)$;
  \item первый момент сохранения газа~\eqref{eq:moment1_gas}
        $\longrightarrow$ определяет $m(t)$.
\end{enumerate}

Таким образом, показатели формы~$n(t)$ и~$m(t)$ эволюционируют
вместе с остальными параметрами и \textbf{перестают быть
адаптационными}.

\subsection{Метод коллокации как альтернатива}

Вместо моментных соотношений можно потребовать, чтобы аппроксимация
точно удовлетворяла исходному PDE в выделенной точке --- например,
на забое ($\xi = 0$). Подставляя
\eqref{eq:pressure_profile}--\eqref{eq:saturation_profile}
в~\eqref{eq:oil_conservation} при $x = 0$:
\begin{equation}\label{eq:collocation_oil}
  \phi\,\frac{\partial}{\partial t}
  \left(\frac{1-S_{g,\max}}{B_o(p_{wf})}\right)
  = \frac{k}{x_f^2}\,\frac{\partial}{\partial\xi}
  \left.\left(
    \frac{k_{ro}(S_g)}{\mu_o(p)\,B_o(p)}\,
    \frac{n\,\Delta p}{1}\,(1-\xi)^{n-1}
  \right)\right|_{\xi=0}.
\end{equation}

Правая часть после дифференцирования содержит $n$, $m$ и известные
функции PVT-свойств, что даёт алгебраическое соотношение для
показателей формы. Аналогичная подстановка в уравнение для газа
при~$\xi = 0$ даёт второе соотношение.

\subsection{Сравнение подходов}

\begin{center}
\begin{tabular}{lcc}
  \hline
  & \textbf{Метод моментов} & \textbf{Коллокация} \\
  \hline
  Тип уравнений для $n$, $m$ & ОДУ & алгебраические \\
  Физический смысл & глобальный (форма профиля) & локальный (точка $x=0$) \\
  Точность вдали от забоя & выше & ниже \\
  Вычислительная сложность & выше (доп.\ интегралы) & ниже \\
  \hline
\end{tabular}
\end{center}

На практике часто комбинируют оба подхода: первый момент определяет
один показатель, коллокация --- другой. Это позволяет сбалансировать
точность и трудоёмкость.

\subsection{Замечание о сходимости}

При увеличении числа моментов (нулевой, первый, второй, \ldots)
приближённое решение систематически сходится к точному. Формально
это следует из полноты системы степенных функций и теоремы Вейерштрасса:
любой непрерывный профиль может быть аппроксимирован полиномом
с произвольной точностью.

Однако уже два момента (нулевой + первый) обычно дают погрешность
порядка нескольких процентов, что достаточно для инженерных оценок.
Дальнейшее увеличение порядка требует усложнения аппроксимации
(добавление членов с различными показателями) и редко оправдано.

%=============================================================
\section{Заключение}
%=============================================================

Метод интегральных соотношений Баренблатта позволяет свести задачу
двухфазной фильтрации в модели Black Oil к системе обыкновенных
дифференциальных уравнений. Основные преимущества:
\begin{itemize}
  \item физически прозрачная структура решения с явным отслеживанием
        фронта двухфазной зоны;
  \item возможность аналитического вычисления интегралов благодаря
        степенной аппроксимации профилей;
  \item правильное воспроизведение качественных закономерностей:
        рост газонасыщенности, увеличение газового фактора,
        конечная скорость распространения фронта;
  \item вычислительная простота --- решение системы ОДУ вместо PDE.
\end{itemize}

Показатели формы $n$, $m$ могут быть определены из сравнения
с численным решением полной задачи или из дополнительных
интегральных соотношений (метод моментов).

\end{document}
